\subsection{Non-Factorizable Projected Descriptions and Quantum Correlations}
  \label{subsec:non-injective-projection}

  The projection $\Pi$ is a fixed mapping with intrinsic finite resolution
  and irreducible projection noise: its pixelisation is set by the spectral
  admissibility threshold $\lambda^*$ (Section~\ref{subsec:relational-projection}) and its capacity is
  bounded by the Born--Infeld saturation condition.
  Observable dynamics is not a dynamics of $\chi$, which is atemporally
  fixed, but a sequence of admissible reprojections of the effective state:
  %
  \begin{equation}
    \chi_{\mathrm{eff}}(t) \;\longrightarrow\; \chi_{\mathrm{eff}}(t+1).
  \end{equation}
  %
  Each step generically loses information (non-injectivity), introduces a
  projective deformation (projection noise), and preserves certain global
  invariants (admissibility constraints).
  What appears as separate particles or subsystems at time~$t$ are
  excitations of $\chi_{\mathrm{eff}}(t)$ that $\Pi$ can resolve as distinct.

  Entanglement arises as a resolution event in this reprojection sequence:
  two degrees of freedom $A$ and $B$ of $\chi_{\mathrm{eff}}(t)$ fall below
  the resolution of $\Pi$ in the relevant observable sector, so that
  $\Pi(A) \approx \Pi(B)$ within that sector.
  They share an effective resolution cell.
  No special or pre-selected configuration of $\chi$ is required, and no
  superdeterminism is implied: whether two excitations fall into the same
  cell is determined by the structure of $\chi_{\mathrm{eff}}(t)$ relative
  to the fixed pixelisation of $\Pi$, exactly as two distinct inputs may
  land on the same pixel of a camera whose resolution is fixed.

  This does not reduce $\chi$ to a passive or contentless mirror.
  $\chi$ does not encode observable values, but it carries the full
  structural information determining which relational configurations are
  admissible.
  The global balance invariants---total spin, total charge, winding
  number---that appear as infra-projectable constraints of the fiber
  $\Pi^{-1}(o_{AB})$ originate in this admissibility structure, not in the
  effective state $\chi_{\mathrm{eff}}(t)$ and not in a classically rich
  hidden state.
  $\chi$ is minimal in observable content and maximal in structural
  constraint: it is a coherence space, not a data store.
  Observable dynamics is therefore not a self-reflection of the universe
  through a deforming mirror, but a sequence of constrained reconstructions
  in which $\chi$ determines what can consistently emerge, without itself
  containing what will emerge.

  The deeper reason for this non-factorizability can be identified at the
  level of the spectral structure of the projection.
  A configuration $\omega \in \Omega_\chi$ may carry a global
  invariant $I(\omega)$---total charge, total angular momentum, winding
  number, or more generally any topological raccordement constraint---that
  belongs to a sector of the substrate which is not individually projectable
  within the effective description.
  Such invariants belong to spectral sectors of $L_\chi$ below the
  projection threshold $\lambda^*$ (Section~\ref{subsec:relational-projection}) and therefore cannot appear
  as independent local observables in the effective description: they are not
  representable as a separable pair $I_A(\omega_A) \oplus I_B(\omega_B)$
  projectable independently into $O_A \times O_B$, but exist once, at the
  level of $\chi$, and can only appear in the effective description as a
  property of the joint fiber $\Pi^{-1}(o_{AB})$.
  The relevant pre-image is therefore not $\Pi^{-1}(o_A) \times \Pi^{-1}(o_B)$
  but the non-factorisable global class $\Pi^{-1}(o_{AB})$, whose elements
  are configurations constrained jointly by $I(\omega)$.

  The type of correlation is determined by which balance constraint is
  infra-projectable, not by a different mechanism.
  Direct correlations ($s_A = s_B$) arise when the infra-projectable
  invariant selects decompositions with equal local values;
  anti-correlations ($s_A + s_B = 0$) arise when the invariant forces
  opposite local values.
  In both cases the constraint is carried by the fiber, not by individual
  pre-images.
  Measurement at~$A$ selects a locally stable reprojection
  $\Pi_A(\omega) = o_A^*$; because the fiber $\Pi^{-1}(o_{AB})$ is
  non-factorizable, this selection simultaneously restricts the class of
  configurations admissible for~$B$ to those compatible with $I(\omega)$
  and $o_A^*$, without any signal propagating from $A$ to $B$.
  The apparent ``collapse'' at~$B$ is a restriction of admissible projected
  descriptions imposed by the global non-decomposability of the fiber, not a
  dynamical influence.

  \subsubsection*{Ontological Monism and the Shared Projection Hypothesis}
    \label{subsec:ontological-monism-shared-projection}

    A single connected excitation in~$\chi$ can be represented as several
    spatially separated effective excitations.
    The observed separation is a property of the projected metric
    representation, not a fundamental separation of the underlying entity.

    \paragraph{Shared Fiber Phase and Spin Correlations.}
      \label{subsec:shared-fiber-phase-spin}
      Spin correlations can be read as correlations of a \emph{shared}
      internal degree of freedom of the projection fiber.
      A measurement at one location selects a locally stable reprojection of
      that shared fiber degree of freedom; the correlated statistics at the
      distant location reflect the restricted set of reprojections still
      compatible with the same underlying configuration.

    \begin{figure}[t]
      \centering
      \begin{tikzpicture}[
        font=\small,
        node distance=10mm,
        box/.style={draw, rounded corners, align=center,
        inner sep=6pt},
        arrow/.style={-Latex, thick},
        note/.style={align=left, font=\footnotesize},
        micro/.style={draw, circle, inner sep=1.5pt}
      ]
    \node[micro] (c1) at (-2.4, 1.5) {};
    \node[micro] (c2) at (-1.2, 1.7) {};
    \node[micro] (c3) at ( 0.0, 1.55) {};
    \node[micro] (c4) at ( 1.2, 1.75) {};
    \node[micro] (c5) at ( 2.4, 1.6) {};
    \node[note, above=3mm of c3, align=center]
    {Admissible $\chi$-configurations\\
    (non-decomposable joint fiber)};
    \node[box, below=16mm of c3,
      minimum width=6.8cm] (chieff)
    {$\chi_{\mathrm{eff}}$\\
    \footnotesize constrained reconstruction\\
    \footnotesize (non-factorizable)};
    \node[box, below left=12mm and 14mm of chieff] (obsA)
    {Outcomes in context $A$};
    \node[box, below right=12mm and 14mm of chieff] (obsB)
    {Outcomes in context $B$};
    \draw[arrow] (c1) -- (chieff.north west);
    \draw[arrow] (c2) -- (chieff.north);
    \draw[arrow] (c3) -- (chieff.north);
    \draw[arrow] (c4) -- (chieff.north);
    \draw[arrow] (c5) -- (chieff.north east);
    \node[note, left=2mm of chieff, anchor=east]
    {\footnotesize non-injective\\
    \footnotesize projection $\pi$};
    \draw[arrow] (chieff)
    -- node[left=2mm, note] {$\mathcal{O}_A$} (obsA);
    \draw[arrow] (chieff)
    -- node[right=2mm, note] {$\mathcal{O}_B$} (obsB);
      \end{tikzpicture}
      \caption{Non-factorizable regime.
      The joint fiber $\Pi^{-1}(o_{AB})$ is non-decomposable: no
      independent pre-images can be assigned to subsystems $A$ and $B$.
      The admissibility structure of $\chi$ determines which balance
      invariants are infra-projectable; it does not select a unique
      pre-existing configuration.
      Different measurement contexts define distinct admissible
      reprojections of the same constrained fiber, producing correlated
      outcomes without superluminal influence.}
      \label{fig:entanglement-two-projections}
      \end{figure}
